\chapter{Introduction and internship objectives}
\section{From data analysis to molecular ensembles}
Molecular descriptors form a familiar data-analysis problem: clustering groups rows with similar measurements. Interpretation requires care because consecutive observations may be correlated, simulation bias changes configuration frequencies, and descriptors can hide structural differences.

This was the setting of the Data Analysis L-31 internship at the University of Messina, with Andrea Ciccolo and Valeria Conti Nibali as internship contacts. The programme introduced protein structure, intrinsic disorder, force fields, molecular dynamics and clustering. Practical work included GROMACS, VMD and Python. The final assignment combined analysis of supplied simulation data with an independent contribution using methods from the degree programme.

An intrinsically disordered peptide is described through many configurations rather than a single fixed structure. Disorder does not imply an absence of structural organization \autocite{vanderlee2014}. Clustering summarizes this diversity, but similarity in measured descriptors alone does not establish a physically distinct conformational family.

\section{The Aβ40 case study}
Cadenelli and colleagues studied monomeric Aβ40 and Aβ42 ensembles using well-tempered metadynamics, density-peak clustering, consensus between local clusters and statistical reweighting \autocite{cadenelli2026}. Metadynamics encourages exploration by adding a bias along selected collective variables; the well-tempered formulation controls its development \autocite{barducci2008}. Frequencies observed under this bias are not automatically equilibrium probabilities.

The supplied materials concern Aβ40: a clustering script, a free-energy surface and 300,003 descriptor observations in three replica segments. The four baseline features describe secondary-structure content, radius of gyration and solvent-accessible surface area (SASA). These measurements make configurations manageable as numerical data; they are not atomic coordinates.

\section{Questions and scope of the report}
\begin{protocolsteps}
\item \textbf{Can the analysis be reconstructed?} Make frame identity, scaling, FES interpolation and weighted summaries explicit and testable.
\item \textbf{What do the descriptors reveal?} Examine redundancy and temporal dependence through correlation, PCA and autocorrelation diagnostics.
\item \textbf{Does removing SASA preserve the fitted groups?} Compare Gaussian mixtures across feature sets, component counts and initialization seeds on shared observations. Correlation with radius of gyration motivates this question but does not answer it; components are not identified as physical states.
\end{protocolsteps}

The feature experiment is separate from the paper reproduction. The FES passes grid checks, but its energy units and construction remain unverified; equilibrium populations are not reported. Code, tables and run records make the completed calculations inspectable \autocite{plumed2019,wratten2021}.
